% ============================================================================
% SPANISCH LANGENTWURF-TEMPLATE
% ============================================================================
% Basiert auf: /templates/lesson-plan/langentwurf-template.tex
% Fachfarbe: #c41e3a (Karminrot)
% Spezifikation: /projects/spanisch/kl07/ANLEITUNG-SPANISCH.md
%
% PFLICHT: Tier 1 (vollständig) für alle Doppelstunden
% ============================================================================

\documentclass[11pt,a4paper]{article}

% ============================================================================
% PACKAGES
% ============================================================================
\usepackage[utf8]{inputenc}
\usepackage[T1]{fontenc}
\usepackage[ngerman]{babel}
\usepackage{geometry}
\usepackage{fancyhdr}
\usepackage{lastpage}
\usepackage{xcolor}
\usepackage{tcolorbox}
\usepackage{enumitem}
\usepackage{tabularx}
\usepackage{longtable}
\usepackage{array}
\usepackage{booktabs}
\usepackage{hyperref}
\usepackage{ifthen}
\usepackage{amssymb}

% ============================================================================
% KONFIGURATION
% ============================================================================

% --- TIER (immer 1 für Spanisch) ---
\newcommand{\plantier}{1}

% --- FACH ---
\newcommand{\fachid}{spanisch}
\newcommand{\fach}{Spanisch}
\newcommand{\fachfarbe}{c41e3a}

% ============================================================================
% VARIABLEN - HIER AUSFÜLLEN
% ============================================================================

% --- Metadaten (Platzhalter) ---
\newcommand{\jahrgangsstufe}{7}
\newcommand{\schulart}{Gesamtschule}
\newcommand{\stundenthema}{[THEMA]}                % z.B. "Rally + Begrüßung"
\newcommand{\reihenthema}{[SEQUENZTHEMA]}          % z.B. "Empezamos"
\newcommand{\stundennummer}{[X]}                   % z.B. "1a"
\newcommand{\reihenstunden}{[N]}                   % Gesamtzahl Doppelstunden
\newcommand{\unterrichtsdatum}{XX.XX.XXXX}
\newcommand{\unterrichtszeit}{XX:XX -- XX:XX}
\newcommand{\raum}{XXX}

% --- Lehrkraft ---
\newcommand{\lehrkraft}{N.N.}
\newcommand{\ausbildungsstand}{Lehrkraft}

% --- Lerngruppe ---
\newcommand{\lerngruppengroesse}{24}
\newcommand{\maennlich}{12}
\newcommand{\weiblich}{12}
\newcommand{\divers}{0}

% --- Kompetenzen/Niveau ---
\newcommand{\gerniveau}{A1}                        % A1, A2, B1
\newcommand{\lehrwerk}{Qué pasa 1}
\newcommand{\lehrwerkseite}{[S. X-Y]}

% ============================================================================
% LAYOUT
% ============================================================================
\geometry{left=2.5cm, right=2.5cm, top=2.5cm, bottom=2.5cm}
\definecolor{fach}{HTML}{\fachfarbe}
\definecolor{gray}{HTML}{666666}
\definecolor{lightgray}{HTML}{f5f5f5}
\definecolor{successgreen}{HTML}{2a7a4b}
\definecolor{warningorange}{HTML}{b35c00}

\tcbuselibrary{skins,breakable}

% Farbige Boxen
\newtcolorbox{infobox}[1][]{
    colback=lightgray,
    colframe=fach,
    fonttitle=\bfseries,
    title=#1,
    breakable
}

\newtcolorbox{scorebox}{
    colback=white,
    colframe=fach,
    fonttitle=\bfseries,
    title=Didaktik-Score,
    breakable
}

% Kopf-/Fußzeile
\pagestyle{fancy}
\fancyhf{}
\fancyhead[L]{\small\textcolor{gray}{\fach{} -- Jg. \jahrgangsstufe{} -- \stundenthema}}
\fancyhead[R]{\small\textcolor{gray}{Langentwurf Tier 1}}
\fancyfoot[C]{\small\thepage{} / \pageref{LastPage}}
\renewcommand{\headrulewidth}{0.4pt}
\renewcommand{\footrulewidth}{0pt}

% ============================================================================
% DOKUMENT
% ============================================================================
\begin{document}

% ============================================================================
% DECKBLATT
% ============================================================================
\begin{titlepage}
    \centering
    \vspace*{2cm}

    {\large\textcolor{gray}{\schulart}}\\[2cm]

    {\Huge\textcolor{fach}{\textbf{Langentwurf}}}\\[0.5cm]
    {\large Tier 1 -- Vollständig}\\[2cm]

    {\LARGE\textbf{\stundenthema}}\\[0.5cm]
    {\large im Rahmen der Unterrichtsreihe}\\[0.3cm]
    {\Large\textit{\reihenthema}}\\[2cm]

    \begin{tabular}{rl}
        \textbf{Fach:} & \fach{} (\gerniveau) \\[0.3cm]
        \textbf{Klasse:} & \jahrgangsstufe \\[0.3cm]
        \textbf{Lehrwerk:} & \lehrwerk, \lehrwerkseite \\[0.3cm]
        \textbf{Datum:} & \underline{\hspace{3cm}} \\[0.3cm]
        \textbf{Zeit:} & \underline{\hspace{3cm}} (Doppelstunde) \\[0.3cm]
        \textbf{Raum:} & \underline{\hspace{3cm}} \\[0.3cm]
    \end{tabular}

    \vfill

    {\large Lehrkraft: \lehrkraft}\\[0.3cm]
    {\normalsize\textcolor{gray}{\ausbildungsstand}}\\[1cm]

\end{titlepage}

% ============================================================================
% INHALTSVERZEICHNIS
% ============================================================================
\tableofcontents
\newpage

% ============================================================================
% 1. BEDINGUNGSANALYSE
% ============================================================================
\section{Bedingungsanalyse}

\subsection{Lerngruppe}
\begin{tabular}{ll}
    Kursgröße: & \lerngruppengroesse{} SuS (\maennlich{} m / \weiblich{} w / \divers{} d) \\
    Jahrgangsstufe: & \jahrgangsstufe \\
    Sprachniveau: & \gerniveau{} (GER) \\
    Fremdsprache: & 2. Fremdsprache \\
\end{tabular}

\subsection{Differenzierung (intern)}

\textbf{WICHTIG:} Keine Sternchen auf Schülermaterial!

\begin{tabular}{|l|p{4cm}|p{4cm}|p{4cm}|}
\hline
& \textbf{[B] Basis} & \textbf{[S] Standard} & \textbf{[E] Erweiterung} \\
\hline
\textbf{Ziel} & Berufsreife & Mittlere Reife & Gymnasium \\
\hline
\textbf{Inhalt} & Rezeptiv, Lückentext & Produktion mit Hilfe & Freie Produktion \\
\hline
\textbf{Hilfen} & Wortschatz, Beispiele & Teils vorgegeben & Nur Impulse \\
\hline
\end{tabular}

\subsection{Besonderheiten}
\begin{itemize}
    \item \textbf{LRS:} [Anzahl] SuS (Nachteilsausgleich: Zeitzugabe)
    \item \textbf{DaZ:} [Anzahl] SuS (Wortschatzarbeit verstärken)
    \item \textbf{Leistungsstark:} [Anzahl] SuS (Zusatzaufgaben)
\end{itemize}

% ============================================================================
% 2. SACHANALYSE
% ============================================================================
\section{Sachanalyse}

\subsection{Sprachliche Strukturen}
\begin{itemize}
    \item \textbf{Grammatik:} [z.B. Verbkonjugation ser, estar]
    \item \textbf{Wortschatz:} [z.B. Begrüßungen, Zahlen 1-20]
    \item \textbf{Phonetik:} [z.B. Aussprache von ñ, rr]
\end{itemize}

\subsection{Kontrastive Betrachtung (Deutsch $\leftrightarrow$ Spanisch)}
\begin{itemize}
    \item \textbf{Positive Transfers:} [z.B. Ähnliche Frageintonation]
    \item \textbf{Interferenzen:} [z.B. Subjektpronomen nicht obligatorisch]
    \item \textbf{Falsche Freunde:} [z.B. ...]
\end{itemize}

\subsection{Kulturelle Aspekte}
\begin{itemize}
    \item [z.B. Begrüßungsrituale in Spanien vs. Deutschland]
    \item [z.B. Duzen/Siezen (tú/usted)]
\end{itemize}

% ============================================================================
% 3. DIDAKTISCHE ANALYSE
% ============================================================================
\section{Didaktische Analyse}

\subsection{Stellung in der Sequenz}
Dies ist die \textbf{\stundennummer. Doppelstunde} der Sequenz ``\reihenthema'' (\reihenstunden{} Doppelstunden gesamt).

\subsection{Kompetenzziele (Rahmenplan MV)}
\begin{itemize}
    \item \textbf{Kommunikativ:} [z.B. Hörverstehen, Sprechen]
    \item \textbf{Sprachlich:} [z.B. Grammatik, Wortschatz]
    \item \textbf{Interkulturell:} [z.B. Landeskunde]
    \item \textbf{Methodisch:} [z.B. Wörterbucharbeit]
\end{itemize}

\subsection{Legitimation}
\textbf{Rahmenplan Spanisch MV (Kl. 7-10):}
\begin{itemize}
    \item Kompetenzbereich: ...
    \item Standard: ...
\end{itemize}

% ============================================================================
% 4. STUNDENZIELE
% ============================================================================
\section{Stundenziele}

\subsection{Grobziel}
Die SuS erweitern ihre \textbf{kommunikative Kompetenz} im Bereich [Hören/Sprechen/Lesen/Schreiben], indem sie [konkrete Handlung].

\subsection{Feinziele (operationalisiert)}
\begin{tabular}{|c|p{8cm}|c|c|}
\hline
\textbf{Nr.} & \textbf{Die SuS können...} & \textbf{AFB} & \textbf{Diff.} \\
\hline
FZ1 & die Vokabeln [X] korrekt aussprechen & I & alle \\
\hline
FZ2 & [Struktur X] in Sätzen anwenden & II & alle \\
\hline
FZ3 & einen kurzen Dialog selbstständig führen & II & [S]/[E] \\
\hline
FZ4 & [Situation X] kulturell einordnen & III & [E] \\
\hline
\end{tabular}

% ============================================================================
% 5. METHODISCHE ANALYSE
% ============================================================================
\section{Methodische Analyse}

\subsection{Phasierung (Doppelstunde = 2 × 45 Min.)}
\begin{enumerate}
    \item \textbf{1. Stunde:}
    \begin{itemize}
        \item Einstieg: [z.B. Bildimpuls]
        \item Erarbeitung: [z.B. Hörtext mit Lücken]
        \item Übung: [z.B. Nachsprechen, Partnerarbeit]
    \end{itemize}
    \item \textbf{2. Stunde:}
    \begin{itemize}
        \item Wiederholung: [z.B. Vokabelquiz]
        \item Anwendung: [z.B. Rollenspiel]
        \item Sicherung: [z.B. Arbeitsblatt]
    \end{itemize}
\end{enumerate}

\subsection{Medien}
\begin{itemize}
    \item Tafel/Whiteboard
    \item Lehrwerk: \lehrwerk, \lehrwerkseite
    \item Arbeitsblatt: AB-[XX][y].pdf
    \item Lernpfad: LP-[XX][y].html (optional, digital)
    \item Audio: [Track X]
\end{itemize}

% ============================================================================
% 6. VERLAUFSPLANUNG
% ============================================================================
\section{Verlaufsplanung}

\textbf{1. Stunde (45 Min.)}

\begin{longtable}{|p{1cm}|p{1.5cm}|p{4cm}|p{3cm}|p{2cm}|p{2cm}|}
\hline
\textbf{Zeit} & \textbf{Phase} & \textbf{Lehrerhandeln} & \textbf{Schülerhandeln} & \textbf{SF} & \textbf{Medien} \\
\hline
\endhead

0--5 & Einstieg & Begrüßung, Tages\-programm, Bildimpuls & Vermutungen äußern & Plenum & Tafel \\
\hline

5--15 & Input & Präsentation [X], Hörtext abspielen & Hören, Notieren & Plenum & Audio \\
\hline

15--25 & Übung 1 & Chorisches Sprechen anleiten & Nachsprechen & Plenum & -- \\
\hline

25--35 & Übung 2 & Partnerarbeit anleiten & Dialoge üben & PA & AB \\
\hline

35--40 & Sicherung & Ergebnisse abfragen & Vorstellen & Plenum & -- \\
\hline

40--45 & Über\-leitung & Ausblick 2. Stunde & Aufräumen & Plenum & -- \\
\hline
\end{longtable}

\vspace{1em}

\textbf{2. Stunde (45 Min.)}

\begin{longtable}{|p{1cm}|p{1.5cm}|p{4cm}|p{3cm}|p{2cm}|p{2cm}|}
\hline
\textbf{Zeit} & \textbf{Phase} & \textbf{Lehrerhandeln} & \textbf{Schülerhandeln} & \textbf{SF} & \textbf{Medien} \\
\hline
\endhead

0--5 & Wieder\-holung & Vokabelquiz & Antworten & Plenum & -- \\
\hline

5--20 & Anwen\-dung & Rollenspiel anleiten & Dialoge führen & PA/GA & Karten \\
\hline

20--35 & Festigung & Arbeitsblatt ausgeben, unterstützen & AB bearbeiten & EA & AB \\
\hline

35--40 & Sicherung & Lösungen besprechen & Vergleichen & Plenum & Tafel \\
\hline

40--45 & Abschluss & Zusammenfassung, HA & Notieren & Plenum & -- \\
\hline
\end{longtable}

\textbf{Legende:} EA = Einzelarbeit, PA = Partnerarbeit, GA = Gruppenarbeit, SF = Sozialform

% ============================================================================
% 7. ERWARTETE SCHWIERIGKEITEN
% ============================================================================
\section{Erwartete Schwierigkeiten}

\begin{tabular}{|p{5cm}|p{8cm}|}
\hline
\textbf{Schwierigkeit} & \textbf{Reaktion} \\
\hline
Aussprache [X] & Chorisches Sprechen, Einzelkorrektur \\
\hline
Verwechslung [X/Y] & Kontrastive Übung, Merkhilfe \\
\hline
Zeitproblem & Puffer: Übung 2 kürzen \\
\hline
Heterogenität & Differenzierte AB-Aufgaben ([B]/[S]/[E]) \\
\hline
\end{tabular}

% ============================================================================
% 8. HAUSAUFGABE
% ============================================================================
\section{Hausaufgabe}

\begin{itemize}
    \item Lehrwerk S. [X], Aufgabe [Y]
    \item Vokabeln lernen: [Liste]
    \item Optional: Lernpfad LP-[XX][y].html bearbeiten
\end{itemize}

% ============================================================================
% 9. ANLAGEN
% ============================================================================
\section{Anlagen}

\begin{enumerate}
    \item Arbeitsblatt (AB-[XX][y].pdf)
    \item Musterlösung (ML-[XX][y].pdf)
    \item Lehrerhinweise (LH-[XX][y].pdf)
    \item Lernpfad (LP-[XX][y].html)
    \item Tafelbild / Folie
    \item Audio-Track [X]
\end{enumerate}

% ============================================================================
% 10. DIDAKTIK-SCORE
% ============================================================================
\newpage
\section{Didaktik-Score}

\begin{scorebox}
\textbf{Bewertung der didaktischen Qualität nach 10 Dimensionen}\\
\textbf{Ziel-Score: $\geq$ 9.0 (Premium-Qualität)}

\vspace{1em}
\begin{tabular}{|c|l|c|c|p{4.5cm}|}
\hline
\textbf{\#} & \textbf{Dimension} & \textbf{Gew.} & \textbf{Score} & \textbf{Notiz} \\
\hline
1 & Differenzierung & 15\% & /10 & \\
\hline
2 & Taxonomie (AFB) & 10\% & /10 & \\
\hline
3 & Lernziel-Alignment & 10\% & /10 & \\
\hline
4 & Aktivierung & 10\% & /10 & \\
\hline
5 & Scaffolding & 10\% & /10 & \\
\hline
6 & Feedback-Qualität & 10\% & /10 & \\
\hline
7 & Sprachliche Klarheit & 10\% & /10 & \\
\hline
8 & Motivation \& Relevanz & 10\% & /10 & \\
\hline
9 & Inklusion (UDL) & 10\% & /10 & \\
\hline
10 & Praktikabilität & 5\% & /10 & \\
\hline
\multicolumn{3}{|r|}{\textbf{GESAMT:}} & \textbf{/10} & \\
\hline
\end{tabular}

\vspace{1em}
\textbf{Bewertungsschwellen:}
\begin{itemize}
    \item[$\Box$] \textcolor{successgreen}{$\geq$ 9.0: \textbf{FREIGABE} (Premium-Qualität)}
    \item[$\Box$] \textcolor{warningorange}{8.0--8.9: Iteration empfohlen}
    \item[$\Box$] 6.0--7.9: Überarbeitung erforderlich
    \item[$\Box$] $<$ 6.0: Grundlegende Neukonzeption
\end{itemize}
\end{scorebox}

\end{document}
